\section{Interface da biblioteca e estados de fronteira}
\label{sec:interface}

\subsection{Núcleo aritmético}

Fixe
\[
 J_{k,m}=[m2^{-k},(m+1)2^{-k}],
 \qquad k\in\Nzero,\quad m\in\mathbb Z,
\]
e, para $r=p/q\in\mathbb Q$ em forma reduzida com $q>0$,
\[
 \height(r)=\max\{|p|,q\}.
\]
Pelas \cref{sec:indexing,sec:generator}, a biblioteca $\Library$ está em bijeção com as
quádruplas
\[
 (k,m,y_0,s)\in\Nzero\times\mathbb Z\times\mathbb Q\times\Qstar.
\]
A quádrupla determina, em $J_{k,m}$, o átomo
\begin{equation}
 \ell(x)=y_0+s(x-x_0),
 \qquad x_0=\left(m+\frac12\right)2^{-k}.
 \label{eq:inherited-atom}
\end{equation}
Sua altura é
\begin{equation}
 \height_{\rm at}(k,m,y_0,s)
 =\max\{2,k+1,|m|+1,\height(y_0),\height(s)\}.
 \label{eq:inherited-height}
\end{equation}
Escreva $\Library_H$ para os átomos de altura no máximo $H$.

As interfaces usadas pela seleção são resumidas na
\cref{tab:interfaces}.
\begin{table}[t]
 \centering
 \caption{Interfaces normativas da biblioteca e da seleção.}
 \label{tab:interfaces}
 \small
 \begin{tabularx}{\linewidth}{@{}lX@{}}
  \toprule
  símbolo & função\\
  \midrule
  $\Gamma,\Gamma^\sharp$ & endereçamentos computáveis do mesmo átomo; o
  segundo ordena por altura e fornece o posto usado no desempate\\
  $\Library,\Library_H$ & biblioteca e sua exaustão pelo calibre
  \eqref{eq:inherited-height}\\
  $\Aclass(K)$ & montagens contínuas de átomos em folhas diádicas\\
  $\Representations_{D,H}^{\Library}(K)$ & representações em árvores de
  profundidade relativa no máximo $D$ e altura no máximo $H$\\
  $\Sigma_{D,H}^{\Library}(K)$ & imagens funcionais dessas representações\\
  $E_{N,H}^{\Library}$ & melhor erro com os dois orçamentos impostos à mesma
  representação\\
  \bottomrule
 \end{tabularx}
\end{table}

Fixaremos uma raiz diádica
\[
 K=J_{k_0,m_0}=[\alpha,\beta]
\]
e denotaremos por $\mathcal D_D(K)$ o conjunto de seus descendentes até
profundidade relativa $D$. Uma árvore é sempre binária, ordenada e cheia:
cada nó é folha ou possui os dois filhos diádicos. Seus conjuntos de folhas
particionam $K$.

Partições binárias recursivas são a estrutura geométrica de CART e de variantes
diádicas de árvores de decisão e aproximação
\citep{breiman1984,donoho1997,blanchard2007}. Aqui elas recebem, além da forma
combinatória, estados racionais de fronteira derivados da biblioteca.

\subsection{Pares admissíveis}

Se $J=J_{k,m}$, escreva $h_J=2^{-k}$. Para valores $u,v\in\mathbb Q$, a única
reta com valores $u$ e $v$ nas extremidades de $J$ é
\begin{equation}
 \ell_{J;u,v}(x)
 =u+\frac{v-u}{h_J}(x-m2^{-k}).
 \label{eq:endpoint-line}
\end{equation}
Seu valor central e sua inclinação são
\[
 y_0=\frac{u+v}{2},
 \qquad s=2^k(v-u).
\]

\begin{definition}[pares de fronteira admissíveis]
\label{def:admissible-pairs}
Para $H\ge2$, defina $\Pairs_H(J)$ como o conjunto dos pares
$(u,v)\in\mathbb Q^2$ que satisfazem
\begin{align}
 u&\ne v,\label{eq:pair-nonzero}\\
 \max\!\left\{\height\!\left(\frac{u+v}{2}\right),
 \height\!\left(2^k(v-u)\right)\right\}&\le H,\label{eq:pair-arithmetic}\\
 k+1&\le H,\qquad |m|+1\le H.\label{eq:pair-geometry}
\end{align}
\end{definition}

\begin{proposition}[dicionário exato de fronteiras]
\label{prop:endpoint-dictionary}
O mapa que associa a um átomo de $\Library_H$ com suporte $J$ seus valores
nas extremidades é uma bijeção sobre $\Pairs_H(J)$. Em particular, se
\[
 \mathcal Q_H=\{r\in\mathbb Q:\height(r)\le H\},
 \qquad q_H=|\mathcal Q_H|,
\]
então
\begin{equation}
 |\Pairs_H(J)|=
 \begin{cases}
 q_H(q_H-1),&\text{se \eqref{eq:pair-geometry} vale},\\
 0,&\text{caso contrário}.
 \end{cases}
 \label{eq:pair-count}
\end{equation}
\end{proposition}

\begin{proof}
Se um átomo tem parâmetros $(k,m,y_0,s)$, seus extremos são
\[
 u=y_0-\frac{s h_J}{2},\qquad
 v=y_0+\frac{s h_J}{2}.
\]
Como $s\ne0$, tem-se $u\ne v$. As identidades inversas
\[
 y_0=\frac{u+v}{2},\qquad s=2^k(v-u)
\]
mostram que a altura atômica não excede $H$ exatamente quando
\eqref{eq:pair-arithmetic}--\eqref{eq:pair-geometry} valem. Logo os dois
mapas são inversos.

Quando a geometria é admissível, pode-se escolher livremente
$y_0\in\mathcal Q_H$ e $s\in\mathcal Q_H\setminus\{0\}$. Há
$q_H(q_H-1)$ escolhas e a bijeção impede colisões. Quando a geometria falha,
nenhum átomo sobre $J$ pertence a $\Library_H$.
\end{proof}

\subsection{Alfabeto finito dos estados}

\begin{definition}[alfabeto de fronteira]
\label{def:value-alphabet}
Defina
\begin{equation}
 \Values_{D,H}(K)
 =\bigcup_{J\in\mathcal D_D(K)}
 \{u,v:(u,v)\in\Pairs_H(J)\}.
 \label{eq:value-alphabet}
\end{equation}
\end{definition}

\begin{lemma}[finitude e suficiência]
\label{lem:value-alphabet}
O conjunto $\Values_{D,H}(K)$ é finito. Todo valor assumido em uma fronteira
de folha por qualquer representação de
$\Representations_{D,H}^{\Library}(K)$ pertence a
$\Values_{D,H}(K)$. Se $S_D=2^{D+1}-1$, então
\begin{equation}
 |\Values_{D,H}(K)|
 \le \min\{2S_Dq_H(q_H-1),\,2|\Library_H|\}
 \le \min\{18S_DH^4,\,36H^6\}.
 \label{eq:value-alphabet-bound}
\end{equation}
\end{lemma}

\begin{proof}
Cada folha é um descendente de $K$ até profundidade $D$ e é rotulada por um
átomo de $\Library_H$. A \cref{prop:endpoint-dictionary} coloca seus dois
valores em \eqref{eq:value-alphabet}, provando a suficiência. Há $S_D$ células
na árvore completa e no máximo $q_H(q_H-1)$ pares por célula; contar os dois
extremos dá a primeira cota. Alternativamente, cada valor contado é extremo
de algum átomo de $\Library_H$, o que dá $2|\Library_H|$. As estimativas
$q_H\le3H^2$ e $|\Library_H|\le18H^6$, provadas na \cref{sec:complexity},
concluem.
\end{proof}

Os estados globais poderão, portanto, ser escritos como pares
$(u,v)\in\Values_{D,H}(K)^2$. Estados não realizáveis recebem custo
$+\infty$. Na implementação, apenas estados realizáveis são armazenados.

O \cref{lem:value-alphabet} foi enunciado para fronteiras de folha, que é o
caso usado adiante nas recorrências. Registre-se explicitamente que ele já
cobre os demais estados: em uma montagem contínua enraizada em $J$, o valor
externo esquerdo é o valor de fronteira esquerdo da folha mais à esquerda de
$T_J$, o valor externo direito é o da folha mais à direita, e o valor
intermediário compartilhado por dois filhos é simultaneamente fronteira
direita de uma folha e fronteira esquerda de outra. Logo todos os estados que
aparecem nas \cref{sec:fixed-dp,sec:variable-dp} pertencem a
$\Values_{D,H}(K)$.

\subsection{As obstruções de escala e de translação}

Uma consequência de \eqref{eq:pair-geometry} merece destaque, porque governa
tudo o que se pode observar em uma instância concreta e é facilmente
confundida com a restrição de profundidade.

\begin{proposition}[células vazias e refinamento uniforme máximo]
\label{prop:translation-obstruction}
Seja $K=J_{k_0,m_0}$ com $m_0\ge0$ e seja
$J=J_{k_0+r,\,2^rm_0+j}$, $0\le j<2^r$, um descendente na profundidade
relativa $r$. Então $\Pairs_H(J)\ne\varnothing$ se e somente se
\begin{equation}
 k_0+r+1\le H
 \qquad\text{e}\qquad
 2^rm_0+j+1\le H.
 \label{eq:two-obstructions}
\end{equation}
Em particular, o número $N_r$ de células não vazias na profundidade $r$ é
\begin{equation}
 N_r=
 \begin{cases}
 \min\{2^r,\ \max\{0,H-2^rm_0\}\},&k_0+r+1\le H,\\
 0,&k_0+r+1>H.
 \end{cases}
 \label{eq:nonempty-cell-count}
\end{equation}
Além disso, \emph{todas} as $2^r$ células dessa profundidade são não vazias
se e somente se
\begin{equation}
 k_0+r+1\le H
 \qquad\text{e}\qquad
 2^r(m_0+1)\le H,
 \label{eq:uniform-refinement-conditions}
\end{equation}
isto é,
\begin{equation}
 r\le
 \min\!\left\{H-k_0-1,\ \log_2\frac{H}{m_0+1}\right\}.
 \label{eq:uniform-refinement-depth}
\end{equation}
\end{proposition}

\begin{proof}
Os índices dos descendentes de $K$ na profundidade $r$ são exatamente
$2^rm_0+j$ com $0\le j<2^r$. Como $m_0\ge0$, a condição de translação
$|m|+1\le H$ equivale a $j\le H-2^rm_0-1$; quando a condição de escala
$k_0+r+1\le H$ vale, contar os inteiros $j\in\{0,\ldots,2^r-1\}$ que
satisfazem essa desigualdade dá \eqref{eq:nonempty-cell-count}. Se a condição
de escala falha, nenhuma célula admite átomo, o que dá a segunda linha.
Todas as células são não vazias exatamente quando a escala é admissível e o
maior índice, $2^rm_0+2^r-1$, ainda satisfaz a condição de translação. Isso é
\eqref{eq:uniform-refinement-conditions}, equivalente a
\eqref{eq:uniform-refinement-depth}. A conclusão sobre $\Pairs_H$ segue de
\eqref{eq:pair-geometry} e da \cref{prop:endpoint-dictionary}.
\end{proof}

\begin{remark}[assimetria espacial do catálogo]
\label{rem:asymmetry}
As duas obstruções em \eqref{eq:two-obstructions} têm naturezas muito
diferentes. A obstrução de escala permite profundidade
$r\le H-k_0-1$, enquanto a de translação permite refinamento uniforme apenas
até $r\le\log_2\bigl(H/(m_0+1)\bigr)$. A profundidade disponível é o menor
desses dois limiares. Para $k_0,m_0$ fixos e $H$ suficientemente grande, o
segundo cresce apenas de modo logarítmico e acaba sendo o limitante; para
outros regimes, a escala pode ser a condição ativa. Quando a translação
limita antes da escala, ela o faz de maneira \emph{espacialmente assimétrica}:
entre as células da mesma profundidade sobrevivem as de menor índice, isto é,
as mais próximas da origem no semieixo não negativo.

Um exemplo é usado adiante. Com $K=[0,1]=J_{0,0}$ e $H=4$, a profundidade
$r=2$ é uniformemente refinável, mas em $r=3$ apenas as células de índice
$m\le3$ admitem átomos: as quatro células que cobrem $[\tfrac12,1]$ têm
$m\in\{4,5,6,7\}$ e satisfazem $\Pairs_4(J)=\varnothing$. O catálogo com
$(D,H)=(3,4)$ refina $[0,\tfrac12]$ até a profundidade $3$ e
$[\tfrac12,1]$ apenas até a profundidade $2$.

Pela \eqref{eq:uniform-refinement-conditions}, uma malha uniforme de passo
$2^{-r}$ em $[0,1]$ está disponível para
$H\ge\max\{2,2^r\}$. Mas registra o preço: para essa malha o orçamento de
altura precisa crescer \emph{linearmente em} $1/h$, e não apenas com
$\log(1/h)$. A \cref{sec:dyadic-asymptotic} separa a constante geométrica sem
truncagem da agenda necessária para realizá-la; as
\cref{sec:approximation-classes,sec:direct-height} incorporam esse acoplamento
aos estratos de aproximação.
\end{remark}
